\section{Conclusion}
\label{sec:conclusion}
This paper presented \adapt{}, an open pipeline that converts a flood-annotated map raster and current weather into a QGroundControl WPL~110 UAV relief mission. It combines colour-based flood extraction, a weather-driven obstacle-inflation rule, boundary drop-point selection, nearest-neighbour ordering, per-leg D*~Lite planning and a flat-earth geodetic conversion with explicit validity checks. It offers two approaches. Either one drone visits every drop point, or a fleet of drones is used, in which each drone serves the cluster of drop points around its own launch base under a battery and payload model. Fleet plans whose routes cross are rejected. The work addresses an intersection that the reviewed literature treats in separate pieces. Its contribution is integration and verification rather than any individual algorithm.

On three published flood maps the pipeline produced missions with 16, 33 and 1 deliveries, all of which passed an independent specification-based validator. The geospatial and mission layers were verified against an independent WGS84 oracle: mission rows match their source pixels within $3.6\times10^{-15}$~degrees, 98\,789 randomised round trips agree within $7\times10^{-8}$~px, and 31 of 33 injected defects were detected. With identical inputs, the multi-base mode reproduced its reference implementation's plans exactly. In the fleet runs, 190 of 206 served flood regions went to a single drone, so each drone in effect took its own cluster. Raising the number of bases from one to four increased the drops served from 10 to 18 of 22 (Varanasi) and from 17 to 34 of 41 (Kanpur) under moderate weather. All 43 resulting missions passed the validator, and no routes crossed. Because the assignment does not balance work, however, one Kanpur base kept 15 to 17 drops throughout. The same evaluation exposed limits of the approach. Predicted-flood obstacles reduce the route's exposure to predicted flooding by 47\% and 76\%, but cause all 14 no-path legs, which are flown as straight segments. Segmentation of a thematic map labels 83\% of it as flood. The nearest-neighbour order is up to 44.5\% longer than optimal on random instances. The per-leg D*~Lite is 61--112 times slower than A* because its incremental reuse is unused. Finally, the supplied maps are not georeferenced and their configured scale is wrong by factors of 46--317, so the generated coordinates are not geographically meaningful.

\adapt{} is therefore best understood as a verified software baseline for image-to-mission research, not as a deployable system. The most important next steps are ground-station and SITL validation of both modes, georeferenced inputs, segmentation ground truth, explicit handling of disconnected drop points in the single-UAV mode, and workload-aware, temporally deconflicted multi-UAV allocation.
